\newpage

\section{Conclusiones}

El desarrollo de este trabajo permitió trasladar a una implementación concreta sobre FPGA varios de los conceptos centrales de arquitectura de computadoras estudiados durante la cursada. La construcción de un procesador \textit{RISC-V} segmentado en cinco etapas implicó no sólo diseñar el camino de datos y la lógica de control, sino también resolver problemas reales de integración entre módulos, temporización, manejo de riesgos y validación funcional, que exceden el análisis puramente teórico de la arquitectura.

Uno de los aspectos más relevantes del proyecto fue la incorporación de una \textit{Debug Unit} como mecanismo de control y observabilidad del sistema. Esta unidad permitió desacoplar la lógica de ejecución del procesador de las tareas de depuración, posibilitando la carga de programas por UART, la ejecución continua o paso a paso, la detención controlada y el volcado del estado interno del sistema. Contar con esta herramienta resultó fundamental para validar el comportamiento del procesador, detectar errores de manera sistemática y facilitar el análisis del tránsito de instrucciones a través del \textit{pipeline}.

Desde el punto de vista funcional, la implementación de mecanismos de \textit{forwarding}, \textit{stall} y \textit{flush} permitió resolver adecuadamente los principales riesgos asociados a una arquitectura segmentada. Las pruebas realizadas mostraron que el procesador mantiene la correctitud de ejecución frente a dependencias de datos, conflictos del tipo \textit{load-use} y cambios de flujo producidos por saltos y bifurcaciones. Del mismo modo, el tratamiento del \texttt{HALT} y el drenado del \textit{pipeline} evolucionaron hacia un esquema más robusto, basado en la detección efectiva del vaciado de los registros inter-etapa, lo que permitió garantizar una finalización consistente de la ejecución.

A lo largo del desarrollo surgieron además distintas dificultades prácticas que enriquecieron el alcance formativo del trabajo. Entre ellas se destacaron el ajuste de la latencia efectiva de las memorias para alinearla con lo esperado por el \textit{pipeline}, la diferenciación conceptual entre pausa externa y terminación interna, y la validación del protocolo UART en escenarios de comunicación más extensos, como los volcados de registros, latches y memoria. La resolución de estos problemas llevó a un diseño más claro, modular y coherente, y obligó a adoptar criterios de validación más precisos tanto en simulación como en hardware real.

En conjunto, los resultados obtenidos indican que el sistema desarrollado cumple con los objetivos propuestos y constituye una implementación sólida y funcional del procesador requerido. Más allá del cumplimiento del trabajo práctico, el proyecto integró conocimientos de diseño digital, arquitectura de computadoras, verificación, temporización e interacción entre hardware y software, ofreciendo una experiencia de desarrollo completa sobre FPGA.